\documentclass[11pt,reqno]{amsart}
\usepackage[a4paper,margin=26mm]{geometry}
\usepackage[T1]{fontenc}
\usepackage[utf8]{inputenc}
\usepackage{lmodern,amsmath,amssymb,amsthm,mathtools,microtype,booktabs}
\usepackage[hidelinks]{hyperref}
\newtheorem{theorem}{Theorem}[section]
\newtheorem{lemma}[theorem]{Lemma}
\newtheorem{proposition}[theorem]{Proposition}
\newtheorem{corollary}[theorem]{Corollary}
\theoremstyle{remark}\newtheorem{remark}[theorem]{Remark}
\newcommand{\Sph}{\mathbb S^2}
\newcommand{\R}{\mathbb R}
\newcommand{\Sym}{\operatorname{Sym}}
\newcommand{\tr}{\operatorname{tr}}
\newcommand{\norm}[1]{\lVert#1\rVert}
\newcommand{\Eh}{\widehat E}
\numberwithin{equation}{section}
\setlength{\emergencystretch}{2em}
\raggedbottom
\title[Eight points: negative exponents]{A compact negative-exponent extension of the eight-point certificate}
\author{}
\date{September 21, 2026}
\begin{document}
\begin{abstract}
We extend the finite three-point certificate in the supplied nonnegative-exponent
proof to the interval $-2<s\le0$. Subtracting an exactly summable centroid term
and dividing by $(s+2)/2$ removes the degeneracy at $s=-2$. The resulting
auxiliary potentials are regular on a closed parameter interval, which is
covered by five degree-$12$ certificates with Taylor order $10$. The extension
uses the same matrix blocks, coefficient correction and equality argument as
the supplied proof, with new compact witnesses and new uniform verifications.
The statement at $s=-2$ itself remains nonunique.
\end{abstract}
\maketitle

\section{Statement and normalization}
For $Y=(y_1,\ldots,y_8)\in(\Sph)^8$, set
\[
 \Eh_s(Y)=\sum_{i<j}\phi_s(y_i\cdot y_j),\qquad
 \phi_s(t)=\begin{cases}
 ((2-2t)^{-s/2}-1)/s,&s\ne0,\\
 -\tfrac12\log(2-2t),&s=0.
 \end{cases}
\]
For $s<0$ the value at $t=1$ is defined by continuity; for $s=0$ it is
$+\infty$. Thus collisions are permitted in the domain. At nonzero $s$, this
energy differs from $s^{-1}\sum_{i<j}|y_i-y_j|^{-s}$ by the constant $-28/s$.
Use the same square-antiprism family and squared-height parameter as in the
supplied manuscript~\cite{source}:
\begin{equation}\label{eq:X}
 X(a)=\left\{\left(\sqrt{1-a}\cos\frac{j\pi}{4},
 \sqrt{1-a}\sin\frac{j\pi}{4},(-1)^j\sqrt a\right):0\le j<8\right\}.
\end{equation}

\begin{theorem}\label{thm:negative}
For each $-2<s\le0$ there is a unique squared height $a_s\in(0,1)$ minimizing
$\Eh_s(X(a))$. For every $Y\in(\Sph)^8$,
\[
 \Eh_s(Y)\ge\Eh_s(X(a_s)),
\]
with equality exactly up to an orthogonal transformation and relabeling.
The finite obligations in the proof below are computer-assisted.
\end{theorem}

Combined with the nonnegative-exponent theorem of the supplied package,
Theorem~\ref{thm:negative} gives the normalized statement for every $s>-2$.
This combination uses the supplied positive-exponent proof; the new computation
does not rerun its 111 finite cells or its large-exponent checks.

\section{Removing the endpoint degeneracy}
Put $q=(s+2)/2$. For $0\le q\le1$ and $-1\le t<1$, define
\begin{equation}\label{eq:integral}
 \Psi_q(t)=2\int_{-1}^t(t-u)(2-2u)^{-q-1}\,du.
\end{equation}
This integral and every parameter derivative are well defined on compact
subsets with $t<1$: the integrand is smooth and its derivatives are uniformly
bounded on the compact integration region. Differentiation under the integral
therefore follows from dominated convergence. In particular,
\begin{equation}\label{eq:derivatives}
 \Psi_q(-1)=\Psi_q'(-1)=0,\qquad
 \Psi_q^{(j)}(t)=2^{j-1}(q+1)_{j-2}(2-2t)^{1-q-j}>0\quad(j\ge2),
\end{equation}
where $(q+1)_0=1$. Twice integrating the equality
$\phi_{2q-2}''(t)=q\Psi_q''(t)$ gives, for $q>0$,
\begin{equation}\label{eq:decompose}
 \phi_{2q-2}(t)=\phi_{2q-2}(-1)+4^{-q}(t+1)+q\Psi_q(t).
\end{equation}
The formula includes $q=1$ through the logarithmic definition. The regular
endpoint potentials are
\begin{align}
 \Psi_0(t)&=(1-t)\log\frac{1-t}{2}+t+1,\label{eq:end0}\\
 \Psi_1(t)&=-\frac12\log\frac{1-t}{2}-\frac{t+1}{4}.\label{eq:end1}
\end{align}
Let $V_q(Y)=\sum_{i<j}\Psi_q(y_i\cdot y_j)$, using the continuous value at
collisions for $q<1$ and $+\infty$ for $q=1$. Every $X(a)$ has centroid zero.
Since
\[
 \sum_{i<j}y_i\cdot y_j=\frac12\left(\norm{\sum_i y_i}^2-8\right),
\]
equation~\eqref{eq:decompose} gives the exact identity
\begin{equation}\label{eq:centroid}
 \Eh_{2q-2}(Y)-\Eh_{2q-2}(X(a))
 =\frac{4^{-q}}2\norm{\sum_i y_i}^2
   +q\bigl(V_q(Y)-V_q(X(a))\bigr)\qquad(0<q\le1).
\end{equation}
At $q<1$ it extends to collisions by continuity. At $q=1$ a colliding
configuration already has infinite energy.

\begin{proposition}[The certified auxiliary statement]\label{prop:aux}
For every $q\in[0,1]$, the energy $V_q$ is uniquely minimized, up to congruence,
by an antiprism $X(a(q))$, with $0<a(q)<1/3$.
\end{proposition}

Theorem~\ref{thm:negative} follows from this proposition and
\eqref{eq:centroid}, with $a_s=a((s+2)/2)$. The multiplier $q$ is strictly
positive in the theorem. At $s=-2$ it vanishes: indeed
$\Eh_{-2}(Y)=-18+\tfrac12\norm{\sum_i y_i}^2$, so every centroid-zero
configuration minimizes that energy. The stronger auxiliary statement at
$q=0$ must not be confused with uniqueness for the original energy at $s=-2$.

\section{The regularized root and scalar minorant}
Put $b_-=2+\sqrt2$, $b_+=2-\sqrt2$, and
\[
 d_A=2(1-a),\quad d_B=4(1-a),\quad
 d_C=2-\sqrt2+b_-a,\quad d_D=2+\sqrt2+b_+a.
\]
These squared distances have unordered-pair multiplicities $8,4,8,8$.
With weights $w=(16,16,-8b_-,-8b_+)$, which sum to zero, define
\begin{equation}\label{eq:Fbar}
 \overline F(q,a)=\begin{cases}
 \displaystyle\sum_{i=A,B,C,D}w_i\frac{e^{-q\log d_i}-1}{q},&q\ne0,\\[2mm]
 \displaystyle-\sum_{i=A,B,C,D}w_i\log d_i,&q=0.
 \end{cases}
\end{equation}
The removable extension is smooth for $0<a<1$, and direct differentiation gives
\begin{align}
 \partial_a V_q(X(a))&=\tfrac12\overline F(q,a),\label{eq:rootderivative}\\
 \partial_a\overline F(q,a)
 &=32d_A^{-q-1}+64d_B^{-q-1}
   +8b_-^2d_C^{-q-1}+8b_+^2d_D^{-q-1}>0.\label{eq:Fbara}
\end{align}
The closed interval certificates below validate a root of $\overline F=0$
with $0<a(q)<1/3$ for every $q\in[0,1]$. Strict positivity in
\eqref{eq:Fbara} makes that root unique on $(0,1)$ and, by
\eqref{eq:rootderivative}, makes it the unique family minimizer. The implicit
function theorem applies, including in a neighborhood of either endpoint.

For clarity, root validation is not based on trusting a numerical root.
The checker constructs a polynomial $p(\tau)$ for $q=c+h\tau$, $|\tau|\le1$,
bounds $|\overline F(q,p(\tau))|$ by $\epsilon$, and bounds
$\partial_a\overline F$ from below by $m>0$ on a padded tube. It checks that
$\epsilon/m$ is smaller than the padding. The mean value theorem and the
intermediate value theorem give the unique exact root in that tube and
$|a(q)-p(\tau)|\le\epsilon/m$. This error is retained in all subsequent models.

Define the contact inner products
\begin{equation}\label{eq:contacts}
 A=a,\quad B=2a-1,\quad C=-a+(1-a)/\sqrt2,\quad D=-a-(1-a)/\sqrt2.
\end{equation}
The checker validates that they lie in $(-1,1)$, are distinct and nonzero,
and contain no opposite pair. Let $H_q$ be the polynomial of degree at most
$12$ agreeing with $\Psi_q$ to multiplicity two at each of $A,B,C,D$, and
to multiplicity five at $-1$. Existence and uniqueness follow from the
invertible confluent interpolation problem at these distinct nodes.
For $t<1$, the Hermite remainder and \eqref{eq:derivatives} give
\begin{equation}\label{eq:minorant}
 \Psi_q(t)-H_q(t)\ge
 \frac{\Psi_q^{(13)}(-1)}{13!}(t+1)^5
       \prod_{\alpha\in\{A,B,C,D\}}(t-\alpha)^2\ge0.
\end{equation}
Here the derivative at the remainder point is at least its value at $-1$,
since $\Psi_q^{(14)}>0$. Equality occurs only at the five interpolation nodes.
For $q<1$, $\Psi_q(1)=2\,4^{-q}/(1-q)$, and taking $t\uparrow1$ in
\eqref{eq:minorant} proves a strictly positive gap also at $t=1$.
For $q=1$ the potential diverges there. Thus collisions cannot attain equality.

\section{What the finite certificates prove}
We use the three-point construction of Cohn and Woo~\cite{CW} in the explicit
form of Section 2 of the supplied manuscript~\cite{source}. We spell out why
its finite criterion applies to the new potentials. Set
$e_q=V_q(X(a(q)))$ and
\[
 P_q(u,v,t)=\frac{H_q(u)+H_q(v)+H_q(t)}3-\frac{e_q}{28},\qquad
 \Delta=1+2uvt-u^2-v^2-t^2.
\]
For $0\le k\le5$, let $Q_0=1$, $Q_1=t-uv$, and
\[
 Q_k=2(t-uv)Q_{k-1}-(1-u^2)(1-v^2)Q_{k-2}.
\]
For $0\le i,j\le6-k$, set $S_{k,ij}=\Sym(u^iv^jQ_k)$, where $\Sym$ is
the average over all six coordinate permutations, and define
\[
 R_k=6S_k+S_k(u,u,1)+S_k(v,v,1)+S_k(t,t,1)+\frac{\delta_{k0}}7\mathbf J.
\]
The full moment $M_k(Y)$ is the sum of $S_k$ over all ordered triples of
labels, allowing repeated labels. For a fixed first point, projecting the
other two points to its tangent plane expresses the double sum through
cosine and sine feature outer products. This proves $M_k(Y)\succeq0$;
vanishing tangent projections are handled by continuity of the polynomial
formula. Separating repeated labels gives
\begin{equation}\label{eq:count}
 \sum_{\text{ordered distinct labels}}R_k=6M_k(Y).
\end{equation}
These are geometric identities, independent of the potential or of $q$.

At the antiprism each $M_k$ has rank at most one: vertex transitivity and
the reflection fixing a vertex make its cosine feature vector common and
its sine vector zero. A verified positive diagonal pivot proves rank one.
Its kernel basis is formed from the columns
$e_f-M_k[r,f]e_r/M_k[r,r]$ and a fixed witness change of basis, giving $N_k$.
For each of the raw monomial vectors of degrees at most $6,5,4$, impose
vanishing on the 21 ordered candidate contact triples. Inverting a checked
$21$-column evaluation minor yields a kernel basis; a fixed witness change
of basis gives the feature vectors $z_0,z_1,z_\Delta$. The candidate's five
unordered types are
\[
 (D,D,A),\ (D,B,C),\ (D,C,A),\ (B,A,A),\ (C,C,A),
\]
with respective multiplicities $8,16,16,8,8$.

The new witnesses validate the coefficient identity
\begin{align}\label{eq:certificate}
 P_q={}&\sum_{k=0}^5\tr(N_kB_kN_k^TR_k)
       +\Sym(z_0^TW_0z_0)\notag\\
     &+\Sym((1-u^2)z_1^TW_1z_1)
       +\Delta\Sym(z_\Delta^TW_\Delta z_\Delta),
\end{align}
where all nine reduced matrices are strictly positive definite.
Their orders are
\[
 (6,5,4,3,2,1;63,35,14).
\]
There are $102$ symmetric polynomial coefficient coordinates and $2807$
symmetric-matrix coordinates. The latter use diagonal entries unchanged
and off-diagonal entries multiplied by $\sqrt2$.

\subsection{Exactness, not a small numerical residual}
At the fixed root, define two linear functionals on symmetric polynomials:
$L_0(f)$ is its ordered-distinct-triple sum on $X(a)$, and $L_1(f)$ is the
height derivative of that sum, with the polynomial coefficients frozen.
They are independent, since
\[
 L_0(1)=336,\quad L_1(1)=0,\quad
 L_1(u^2+v^2+t^2)=1152(3a-1)\ne0.
\]
Every certificate column belongs to $\ker L_0\cap\ker L_1$. For SOS columns
this follows from contact vanishing to second order in the product. For
moment columns, the frozen compression of the positive semidefinite full
moment is zero at the candidate; each quadratic form has a minimum there,
so its first derivative is zero, and polarization gives the matrix claim.
The target belongs to the same common kernel, because interpolation gives
$L_0(P_q)=0$, while first-derivative interpolation and
\eqref{eq:rootderivative} give $L_1(P_q)=6\overline F(q,a)=0$.
This is the only stationarity input needed here; it holds also at $q=0$.

Let $A(q)w=b(q)$ be the full coefficient system. The witness chooses $100$
rows $R$, a fixed correction matrix $C_r$, a left preconditioner $L$ and
an initial vector $w_0$. Set
\[
 T=LA_RC_r,\qquad r=L(b_R-A_Rw_0),\qquad
 w=w_0+C_rT^{-1}r.
\]
The exact checker proves $T$ invertible uniformly. Thus $A_R$ has rank
$100$, its image fills the $100$-dimensional common kernel, and restriction
to the selected rows is injective on that kernel. The selected residual
is exactly zero by construction; therefore the full residual is exactly
zero. The checker encloses this corrected $w$, including its complete
Taylor and inverse errors, and validates all its reduced matrices as
positive definite. A small floating-point residual is never substituted
for this argument.

\subsection{Global inequality and equality}
Summing \eqref{eq:certificate} over ordered distinct labels gives
$12(E_{H_q}(Y)-e_q)\ge0$: the moment terms are nonnegative by
\eqref{eq:count}, and the remaining terms are nonnegative on spherical
triples because $1-u^2\ge0$ and $\Delta\ge0$. Combining this with
\eqref{eq:minorant} proves $V_q(Y)\ge e_q$ for every configuration, not
merely for symmetric or stationary competitors.

For equality, the scalar gap first excludes collisions and restricts
inner products to $\{-1,A,B,C,D\}$. An antipodal pair is impossible: a
third point would give an opposite pair of contact values, or zero, both
excluded by the node tests. Strict positivity of $B_0$, and invertibility
of its fixed change of basis, force each per-vertex moment vector to equal
the candidate's moment vector: it is proportional to that vector by the
kernel condition, and its zeroth coordinate is $8$. The Vandermonde system
at the four distinct nodes then gives counts $n_A=2,n_B=1,n_C=n_D=2$ at
every vertex. The first moments also give centroid zero.

Strict positivity of $W_0$ forces $z_0$ to vanish at every equality triple.
The checker validates a nonzero feature at $(A,A,A),(A,A,C),(A,A,D)$,
so two incident $A$ edges force the opposite edge to be a $B$ edge.
The graph of $A$ edges is two-regular, has no triangles, and cannot have
a cycle of length at least five (it would give two distinct $B$ neighbors
at each vertex). It is consequently two four-cycles. Each cycle has
adjacent inner product $a$ and opposite inner product $2a-1$, and therefore
is a Euclidean square of center squared norm $a$ and radius
$\sqrt{1-a}$. Centroid zero makes the two centers opposite. The remaining
inner products $C,D$ force planar angle cosines $\pm1/\sqrt2$, hence an
offset of $\pi/4$ modulo $\pi/2$. The configuration is exactly
\eqref{eq:X}, up to congruence. This proves Proposition~\ref{prop:aux}
once the finite obligations in the next section are validated.

\section{Endpoint-safe evaluation and verified cells}
The inherited interval arithmetic and Taylor operations are unchanged.
The new evaluator avoids dividing by $q$ or $q-1$ on a cell containing an
endpoint. Define the entire functions
\[
 E(z)=\sum_{n\ge0}\frac{z^n}{(n+1)!},\qquad
 D(a,b)=\sum_{n\ge0}\frac{\sum_{j=0}^n a^j b^{n-j}}{(n+2)!}.
\]
Here $D(a,b)=[0,a,b]\exp$ is the second divided difference, with its
removable values at repeated nodes. If $\ell=\log(4/(2-2t))$, then
\begin{align}
 \Psi_q(t)&=2\,4^{-q}\ell^2D(-\ell,(q-1)\ell),\label{eq:dd}\\
 \Psi_q'(t)&=4^{-q}\ell E(q\ell),\qquad
 \overline F(q,a)=-\sum_iw_i\log(d_i)E(-q\log d_i).
\end{align}
For distinct parameters formula~\eqref{eq:dd} follows from the explicit
divided-difference formula and \eqref{eq:decompose}; analyticity then
extends it to repeated nodes, including $q=0,1$. Alternatively both
expressions have second $t$ derivative \eqref{eq:derivatives} and vanish
to first order at $-1$. The higher derivatives are evaluated directly
by \eqref{eq:derivatives}.

Both series are absolutely and uniformly convergent on bounded sets.
Truncating after $n=M$, with $|z|\le R$ or $|a|,|b|\le R$, gives in either
case the valid error bound
\begin{equation}\label{eq:tail}
 \frac{e^R R^{M+1}}{(M+2)!}.
\end{equation}
For $E$, factor the first omitted power and bound the remaining factorial
ratio by $1/k!$. For $D$, bound its numerator by $(n+1)R^n$ and use
$(n+1)/(n+2)!\le1/(n+1)!$ before the same estimate. The implementation
uses $M=32$ on these order-$10$ cells. Arithmetic Taylor remainders and
the analytic series tails are both retained; they are not interchangeable.

The parameter cells, read as exact rational numbers, are
\[
 [0,1/5],\ [1/5,2/5],\ [2/5,3/5],\ [3/5,4/5],\ [4/5,1].
\]
Each uses degree $12$, contact multiplicities $(2,2,2,2)$, endpoint
multiplicity $5$, Taylor order $10$, and certificate multiplier $1$.
The witnesses are new numerical proposals; the exact checker accepts them
only after rebuilding the coefficient system and all bounds uniformly.
\begin{table}[htbp]
\centering
\caption{Fresh uniform verification: upward-rounded bounds.}
\begin{tabular}{@{}ccc@{}}
\toprule
Auxiliary $q$ interval & Maximum PSD error $\delta$ & Inverse residual \\
\midrule
$[0,1/5]$ & 0.613964 & $5.436\times10^{-9}$ \\
$[1/5,2/5]$ & 0.619303 & $5.593\times10^{-9}$ \\
$[2/5,3/5]$ & 0.534026 & $4.266\times10^{-9}$ \\
$[3/5,4/5]$ & 0.438488 & $4.210\times10^{-9}$ \\
$[4/5,1]$ & 0.319059 & $3.980\times10^{-9}$ \\
\bottomrule
\end{tabular}
\end{table}

In the table, $\delta$ bounds the norm of each preconditioned reduced
matrix minus the identity; the reported maximum is over all nine blocks.
Thus $\delta<1$ proves strict positive definiteness. The inverse residual
bounds the complete polynomial inverse residual and must also be less
than one. Displayed decimal bounds are rounded upward; full outward
machine bounds, root tubes, basis checks and equality exclusions are
retained in the individual \path{verified.json} files. Exact rational
aggregation of the five closed cells gives $[0,1]$ with no gap.

\subsection{Scope and reproducibility}
The publication records contain five fresh uniform cell computations,
followed by an exact rational coverage check. Source and witness hashes
were recorded before each cell and checked unchanged afterward. Hashes
bind inputs; they do not prove inequalities. Numerical roots, inverses,
eigenbases and matrix reconstructions from NumPy or mpmath are untrusted
proposals. Acceptance uses integer interval arithmetic, integer matrix
products, integer square roots and explicit rational series tails.

The local verification used the included GMP exact-integer matrix adapter
in place of python-flint. Its matrix operations are arbitrary-precision
integer construction, multiplication, addition and export; it introduces
no floating-point arithmetic. This backend change is disclosed in every
run record. The default replay instead uses python-flint; either backend
must satisfy the same exact-integer interface. The retained positive
package is supplied only as the mathematical source of the finite
criterion and the complementary nonnegative theorem. The new negative
checks do not depend on reading its old success reports.

This is a computer-assisted proof with an explicit mathematical reduction,
not a Lean formalization. It relies on the correctness of the mathematical
argument, checker indexing and control flow, Python, the integer backend,
and the execution platform. To recompute the whole extension using a
fresh output directory, run
\begin{verbatim}
python -B cap.py replay --jobs 4 --output new-replay
\end{verbatim}
For the optional GMP backend, append \verb|--backend gmp|; a C++17 compiler
and GMP development libraries are required. A complete replay recomputes
all five cells; an integrity check or an isolated successful cell is not
a complete replay. See \path{README.md} for the precise command and
publication-record provenance.

\begin{thebibliography}{9}
\bibitem{source} Supplied manuscript, \emph{Optimal eight-point configurations
for all nonnegative exponents}, dated September 18, 2026, especially
Section 2. The unmodified source is retained as
\path{sources/positive_manuscript.tex}; its original archive location is
\path{Thomson_8_Points-main/CAP_General_s/manuscript/paper.tex}.
\bibitem{CW} H. Cohn and J. Woo, \emph{Three-point bounds for energy
minimization}, J. Amer. Math. Soc. \textbf{25} (2012), 929--958.
\end{thebibliography}
\end{document}
